% ESTILO: ISO 690, autor-año (biblatex-iso690)
% cita-iso690-config.tex
\usepackage[
  style=iso-authoryear,
  sortcites=true,
  citereset=chapter,
  backend=biber,
  backref=false,
  autolang=hyphen,
  indexing=true
]{biblatex}

% ============================================================
% AUTORÍA EN VERSALITAS, SOLO EN LA BIBLIOGRAFÍA
% ------------------------------------------------------------
% biblatex-iso690 compone los apellidos en caja alta por
% omisión: MOYANO, Alberto. Se reemplaza por versalitas parejas
% en los cuatro componentes del nombre, que es lo que produce
% philosophy-modern con scauthorsbib y lowscauthors juntas.
%
% EL PLEGADO A MINÚSCULA ES IMPRESCINDIBLE: \textsc SOBRE UNA
% INICIAL ALTA LA DEJA ALTA, Y DA «Moyano» EN LUGAR DE «moyano».
% MEDIDO: \textsc{Moyano} OCUPA 39,97 pt Y \textsc{moyano} 37,56,
% PORQUE LA M SIGUE SIENDO MAYÚSCULA EN EL PRIMERO.
%
% EL APELLIDO VA POR \familynameformat, QUE EL PAQUETE DEFINE
% COMO PUNTO DE EXTENSIÓN Y APLICA CON \AtBeginBibliography.
%
% LOS OTROS TRES VAN POR MACROS DE biblatex QUE APLICAN EN TODAS
% PARTES. POR ESO SE REDEFINEN DENTRO DE \AtBeginBibliography:
% ASÍ EL CAMBIO QUEDA ACOTADO A LA BIBLIOGRAFÍA Y LAS CITAS DEL
% TEXTO SIGUEN EN REDONDA.
%
% LAS CONJUNCIONES —LA «y» ENTRE AUTORES, EL «et al.»— NO PASAN
% POR NINGUNA DE ESTAS MACROS: SON SEPARADORES QUE PONE biblatex
% Y QUEDAN EN REDONDA, QUE ES LO CORRECTO.
% ============================================================
\renewcommand{\familynameformat}[1]{\textsc{\MakeLowercase{#1}}}

\AtBeginBibliography{%
	\renewcommand*{\mkbibnamegiven}[1]{\textsc{\MakeLowercase{#1}}}%
	\renewcommand*{\mkbibnameprefix}[1]{\textsc{\MakeLowercase{#1}}}%
	\renewcommand*{\mkbibnamesuffix}[1]{\textsc{\MakeLowercase{#1}}}%
}

\renewbibmacro*{bibindex}{%
	\ifbibindex
	{\indexnames{author}\indexnames{editor}\indexnames{translator}}
	{}}
\renewbibmacro*{citeindex}{%
	\ifciteindex
	{\indexnames{labelname}}
	{}}

% --- Sangría de referencias ---
\setlength{\bibhang}{2em}

% --- Definiciones para el tratamiento de las url ---
\defcounter{biburlnumpenalty}{3000}
\defcounter{biburlucpenalty}{6000}
\defcounter{biburllcpenalty}{9000}

% REDISEÑO PARA OBTENER MODO LARGO EN ALGUNAS ENTRADAS
% LAS CADENAS QUE biblatex NO TRAE (FUNCIONES DE EDITOR, TIPOS DE
% TESIS, CINE, \sq Y \sqq) ESTÁN EN cita-cadenas-es.tex, COMÚN A
% TODOS LOS ESTILOS. ACÁ SOLO SE REDEFINEN LAS QUE spanish.lbx YA
% TRAE: SON DECISIONES DE ESTE ESTILO
\DefineBibliographyStrings{spanish}{%
	inpreparation    = {en preparación},
	submitted        = {enviado},
	forthcoming      = {de próxima aparición},
	inpress          = {en imprenta},
	prepublished     = {previamente publicado},
	part             = {tomo},
	january          = {enero},
	february         = {febrero},
	march            = {marzo},
	april            = {abril},
	may              = {mayo},
	june             = {junio},
	july             = {julio},
	august           = {agosto},
	september        = {septiembre},
	october          = {octubre},
	november         = {noviembre},
	december         = {diciembre},
	see              = {véase},
	seealso          = {véase también},
	backrefpage      = {re\-fe\-ren\-cia ci\-ta\-da en pági\-na},
	backrefpages     = {re\-fe\-ren\-cia ci\-ta\-da en pági\-nas},
	seenote          = {véase nota},
	quotedin         = {citado en},
	idem             = {ídem},
	idemsf           = {ídem},
	idemsm           = {ídem},
	idemsn           = {ídem},
	idempf           = {ídem},
	idempm           = {ídem},
	idempn           = {ídem},
	idempp           = {ídem},
	ibidem           = {\emph{ibidem}},
	nodate           = {sin fecha},
	withcommentator  = {con comentario de},
	withannotator    = {con notas de},
	withintroduction = {con introducción de},
	withforeword     = {con prólogo de},
	withafterword    = {con epílogo de},
	bycollaborator   = {con colaboración de},
	andothers        = {\emph{et al\adddot}},
	mathesis         = {Tesis de Maestría},
	phdthesis        = {Tesis de Doctorado},
	urlseen          = {visitado el\addspace},
	translationof    = {trad\adddot\addspace de},
	translationas    = {original publicado en},
	url              = {recuperado de},
	newseries        = {nueva época},
	oldseries        = {antigua época},
	byeditorfo       = {ed\adddotspace y pról\adddotspace por},
	byeditorco       = {ed\adddotspace y com\adddotspace por},
	pagetotal        = {págs},
	pagetotals       = {págs},
	techreport       = {informe técnico},
	resreport        = {reporte de investigación},
	file             = {archivo},
	patent           = {patente},
	patentde         = {patente alemana},
	patenteu         = {patente europea},
	patentfr         = {patente francesa},
	patentuk         = {patente británica},
	patentus         = {patente estadounidense},
	patreq           = {solicitud de patente},
	patreqde         = {solicitud de patente alemana},
	patreqeu         = {solicitud de patente europea},
	patreqfr         = {solicitud de patente francesa},
	patrequk         = {solicitud de patente británica},
	patrequs         = {solicitud de patente estadounidense},
	bycompiler       = {comp\adddotspace por},
	issue            = {n\sptext{o}},
	number           = {n\sptext{o}},
	volume           = {vol\adddot},
	volumes          = {vols\adddot},
}

% bytranslatorfo EN SU FORMA LARGA. VA APARTE PORQUE \lbx@lfromlang
% NECESITA \makeatletter: DENTRO DEL BLOQUE GENERAL SE LEÍA \lbx SEGUIDO
% DE TEXTO (EL MISMO ERROR QUE SE CORRIGIÓ EN VERONA). EN ISO 690 NO
% APARECÍA PORQUE EL ESTILO NO USA LA CADENA EN LOS CASOS MEDIDOS
\makeatletter
\DefineBibliographyStrings{spanish}{%
	bytranslatorfo   = {traducido \lbx@lfromlang\ y prologado por},
}
\makeatother
